% GENERATED FILE. Edit canonical lessons, then run npm run generate:books.
% canonical-source-hash: fnv1a64:9fc9ae24198a517e
% canonical-lessons: JA-C164-kayoubi, JA-C164-suiyoubi, JA-C164-mokuyoubi, JA-C164-kinyoubi, JA-C164-nichiyoubi

\chapter{Tuesday Through to Sunday}
\label{ch:ja-tuesday-through-to-sunday}
\hlchaptermodality{164}

\begin{chapteropening}
\textbf{By the end of this chapter:} \emph{I can name the days of the week in my plans.}

Complete the last lesson of chapter 164: I can name the days of the week in my plans.
\end{chapteropening}

\section[\texorpdfstring{\ja{かようび} (kayōbi)}{かようび (kayōbi)}]{\ja{かようび} (kayōbi) — Tuesday}

\label{lesson:JA-C164-kayoubi}



\begin{quote}
Before the new one: say the Japanese for being out, then the Japanese for hygienic.
\end{quote}

\subsection*{You'll want to know: \ja{かようび}}
\textbf{\ja{かようび}} — \emph{kayōbi} — \textquotedblleft{}Tuesday\textquotedblright{}.

\begin{quote}
\textbf{\ja{らいしゅうの} \ja{かようびに} \ja{ともだちに} \ja{あいます。}} — \emph{raishū no kayōbi ni tomodachi ni aimasu} — \textquotedblleft{}I will meet a friend next Tuesday.\textquotedblright{}
\end{quote}

\begin{cousinweb}
\textbf{\ja{かようび}} (\emph{kayōbi}) is a mixed word: Sino-Japanese \emph{ka}, \textquotedblleft{}fire\textquotedblright{} (Middle Chinese \emph{xwaX}), and \emph{yō}, \textquotedblleft{}planet, weekday\textquotedblright{} (\emph{yewH}), borrowed from Chinese, with native \emph{hi}, \textquotedblleft{}day\textquotedblright{}, voiced to \emph{bi} (JMdict, KANJIDIC2; Middle Chinese in Baxter's transcription).
\end{cousinweb}

\begin{grammarlens}[title={What you've built}]
One more word for everyday life at A2.
\end{grammarlens}

\subsection*{Your turn}
\begin{itemize}
\raggedright
  \item \emph{Say it:} \emph{kayōbi}
  \item \emph{Say it:} \emph{kayōbi}, once more
  \item \emph{Say it:} \emph{raishū no kayōbi ni tomodachi ni aimasu} — I will meet a friend next Tuesday.
  \item \emph{Recall:} say the Japanese for being out, then the Japanese for hygienic, then say \emph{kayōbi} again
\end{itemize}

\begin{tcolorbox}[breakable,colback=teal!4,colframe=teal!35!black,title={Before you move on}]
What does \ja{かようび} mean? (\textquotedblleft{}Tuesday\textquotedblright{}.) Say it once more.
\end{tcolorbox}
\section[\texorpdfstring{\ja{すいようび} (suiyōbi)}{すいようび (suiyōbi)}]{\ja{すいようび} (suiyōbi) — Wednesday}

\label{lesson:JA-C164-suiyoubi}



\begin{quote}
Before the new one: say the Japanese for hygienic, then the Japanese for Tuesday.
\end{quote}

\subsection*{You'll want to know: \ja{すいようび}}
\textbf{\ja{すいようび}} — \emph{suiyōbi} — \textquotedblleft{}Wednesday\textquotedblright{}.

\begin{quote}
\textbf{\ja{すいようびは} \ja{あめでしょう。}} — \emph{suiyōbi wa ame deshō} — \textquotedblleft{}It will probably rain on Wednesday.\textquotedblright{}
\end{quote}

\begin{cousinweb}
\textbf{\ja{すいようび}} (\emph{suiyōbi}) is a mixed word: Sino-Japanese \emph{sui}, \textquotedblleft{}water\textquotedblright{} (Middle Chinese \emph{sywijX}), and \emph{yō}, \textquotedblleft{}planet, weekday\textquotedblright{} (\emph{yewH}), borrowed from Chinese, with native \emph{hi}, \textquotedblleft{}day\textquotedblright{}, voiced to \emph{bi} (JMdict, KANJIDIC2; Middle Chinese in Baxter's transcription).
\end{cousinweb}

\begin{grammarlens}[title={What you've built}]
One more word for everyday life at A2.
\end{grammarlens}

\subsection*{Your turn}
\begin{itemize}
\raggedright
  \item \emph{Say it:} \emph{suiyōbi}
  \item \emph{Say it:} \emph{suiyōbi}, once more
  \item \emph{Say it:} \emph{suiyōbi wa ame deshō} — It will probably rain on Wednesday.
  \item \emph{Recall:} say the Japanese for hygienic, then the Japanese for Tuesday, then say \emph{suiyōbi} again
\end{itemize}

\begin{tcolorbox}[breakable,colback=teal!4,colframe=teal!35!black,title={Before you move on}]
What does \ja{すいようび} mean? (\textquotedblleft{}Wednesday\textquotedblright{}.) Say it once more.
\end{tcolorbox}
\section[\texorpdfstring{\ja{もくようび} (mokuyōbi)}{もくようび (mokuyōbi)}]{\ja{もくようび} (mokuyōbi) — Thursday}

\label{lesson:JA-C164-mokuyoubi}



\begin{quote}
Before the new one: say the Japanese for Tuesday, then the Japanese for Wednesday.
\end{quote}

\subsection*{You'll want to know: \ja{もくようび}}
\textbf{\ja{もくようび}} — \emph{mokuyōbi} — \textquotedblleft{}Thursday\textquotedblright{}.

\begin{quote}
\textbf{\ja{もくようびに} \ja{はたらく} \ja{つもりです。}} — \emph{mokuyōbi ni hataraku tsumori desu} — \textquotedblleft{}I plan to work on Thursday.\textquotedblright{}
\end{quote}

\begin{cousinweb}
\textbf{\ja{もくようび}} (\emph{mokuyōbi}) is a mixed word: Sino-Japanese \emph{moku}, \textquotedblleft{}tree, wood\textquotedblright{} (Middle Chinese \emph{muwk}), and \emph{yō}, \textquotedblleft{}planet, weekday\textquotedblright{} (\emph{yewH}), borrowed from Chinese, with native \emph{hi}, \textquotedblleft{}day\textquotedblright{}, voiced to \emph{bi} (JMdict, KANJIDIC2; Middle Chinese in Baxter's transcription).
\end{cousinweb}

\begin{grammarlens}[title={What you've built}]
One more word for everyday life at A2.
\end{grammarlens}

\subsection*{Your turn}
\begin{itemize}
\raggedright
  \item \emph{Say it:} \emph{mokuyōbi}
  \item \emph{Say it:} \emph{mokuyōbi}, once more
  \item \emph{Say it:} \emph{mokuyōbi ni hataraku tsumori desu} — I plan to work on Thursday.
  \item \emph{Recall:} say the Japanese for Tuesday, then the Japanese for Wednesday, then say \emph{mokuyōbi} again
\end{itemize}

\begin{tcolorbox}[breakable,colback=teal!4,colframe=teal!35!black,title={Before you move on}]
What does \ja{もくようび} mean? (\textquotedblleft{}Thursday\textquotedblright{}.) Say it once more.
\end{tcolorbox}
\section[\texorpdfstring{\ja{きんようび} (kin'yōbi)}{きんようび (kin'yōbi)}]{\ja{きんようび} (kin'yōbi) — Friday}

\label{lesson:JA-C164-kinyoubi}



\begin{quote}
Before the new one: say the Japanese for Wednesday, then the Japanese for Thursday.
\end{quote}

\subsection*{You'll want to know: \ja{きんようび}}
\textbf{\ja{きんようび}} — \emph{kin'yōbi} — \textquotedblleft{}Friday\textquotedblright{}.

\begin{quote}
\textbf{\ja{きんようびに} \ja{うちに} \ja{かえります。}} — \emph{kin'yōbi ni uchi ni kaerimasu} — \textquotedblleft{}I will go home on Friday.\textquotedblright{}
\end{quote}

\begin{cousinweb}
\textbf{\ja{きんようび}} (\emph{kin'yōbi}) is a mixed word: Sino-Japanese \emph{kin}, \textquotedblleft{}gold, metal\textquotedblright{} (Middle Chinese \emph{kim}), and \emph{yō}, \textquotedblleft{}planet, weekday\textquotedblright{} (\emph{yewH}), borrowed from Chinese, with native \emph{hi}, \textquotedblleft{}day\textquotedblright{}, voiced to \emph{bi} (JMdict, KANJIDIC2; Middle Chinese in Baxter's transcription).
\end{cousinweb}

\begin{grammarlens}[title={What you've built}]
One more word for everyday life at A2.
\end{grammarlens}

\subsection*{Your turn}
\begin{itemize}
\raggedright
  \item \emph{Say it:} \emph{kin'yōbi}
  \item \emph{Say it:} \emph{kin'yōbi}, once more
  \item \emph{Say it:} \emph{kin'yōbi ni uchi ni kaerimasu} — I will go home on Friday.
  \item \emph{Recall:} say the Japanese for Wednesday, then the Japanese for Thursday, then say \emph{kin'yōbi} again
\end{itemize}

\begin{tcolorbox}[breakable,colback=teal!4,colframe=teal!35!black,title={Before you move on}]
What does \ja{きんようび} mean? (\textquotedblleft{}Friday\textquotedblright{}.) Say it once more.
\end{tcolorbox}
\section[\texorpdfstring{\ja{にちようび} (nichiyōbi)}{にちようび (nichiyōbi)}]{\ja{にちようび} (nichiyōbi) — Sunday}

\label{lesson:JA-C164-nichiyoubi}



\begin{quote}
Before the new one: say the Japanese for Thursday, then the Japanese for Friday.
\end{quote}

\subsection*{You'll want to know: \ja{にちようび}}
\textbf{\ja{にちようび}} — \emph{nichiyōbi} — \textquotedblleft{}Sunday\textquotedblright{}.

\begin{quote}
\textbf{\ja{らいしゅうの} \ja{にちようびは} \ja{うちで} \ja{やすみます。}} — \emph{raishū no nichiyōbi wa uchi de yasumimasu} — \textquotedblleft{}Next Sunday I will rest at home.\textquotedblright{}
\end{quote}

\begin{cousinweb}
\textbf{\ja{にちようび}} (\emph{nichiyōbi}) is a mixed word: Sino-Japanese \emph{nichi}, \textquotedblleft{}sun, day\textquotedblright{} (Middle Chinese \emph{nyit}), and \emph{yō}, \textquotedblleft{}planet, weekday\textquotedblright{} (\emph{yewH}), borrowed from Chinese, with native \emph{hi}, \textquotedblleft{}day\textquotedblright{}, voiced to \emph{bi} (JMdict, KANJIDIC2; Middle Chinese in Baxter's transcription).
\end{cousinweb}

\begin{grammarlens}[title={What you've built}]
One more word for everyday life at A2.
\end{grammarlens}

\subsection*{Your turn}
\begin{itemize}
\raggedright
  \item \emph{Say it:} \emph{nichiyōbi}
  \item \emph{Say it:} \emph{nichiyōbi}, once more
  \item \emph{Say it:} \emph{raishū no nichiyōbi wa uchi de yasumimasu} — Next Sunday I will rest at home.
  \item \emph{Recall:} say the Japanese for Thursday, then the Japanese for Friday, then say \emph{nichiyōbi} again
\end{itemize}

\begin{tcolorbox}[breakable,colback=teal!4,colframe=teal!35!black,title={Before you move on}]
What does \ja{にちようび} mean? (\textquotedblleft{}Sunday\textquotedblright{}.) Say it once more.
\end{tcolorbox}
